\chapter{Conclusion}
\label{ch:conclusion}

\section{Summary of contributions}

This dissertation developed and validated deep-learning methodology for mean-field backward stochastic differential equations (BSDEs), with a focused application to the Cont--Xiong (2024) dealer market-making model. The four contributions, restated with the evidence supporting them:

\paragraph{(C1) Validated neural Bellman solver.}  A stationary neural Bellman solver ([\texttt{solver\_cx.py}](../../solver_cx.py)) for the Cont--Xiong model, trained with a fictitious-play outer loop, achieves spread error of $0.79\%$ and value-function error of $1.28\%$ against the exact Algorithm~1 solution at $N = 2$. The error falls to $0.0001\%$ when the solver is reformulated as a boundary-patched direct-$V$ method. The method scales to larger inventory grids and continuous inventory where the tabular method becomes impractical. Chapter~\ref{ch:bellman}.

\paragraph{(C2) Architectural diagnostics for MV deep BSDE solvers.}  Three independent failure modes that silently suppress mean-field signal in naive deep BSDE solvers were identified, characterised, and fixed:
\begin{itemize}[leftmargin=2em]
  \item \emph{Generator bypass}: the BSDE generator did not depend on the law embedding, severing the gradient path.
  \item \emph{BatchNorm erasure}: broadcast law features have zero batch variance and are annihilated by BatchNorm.
  \item \emph{DeepSets representation collapse}: mean-pooled permutation-invariant encoders cancel symmetric inputs.
\end{itemize}
After correcting all three, the learned competitive factor $h(\Phi(\mu))$ varies by $4\times$ across population shapes and optimal quotes shift by $2\times$; both effects vanish if any single failure mode remains. The findings are robust across two encoder families (moments, quantiles) and generalise beyond the CX model. Chapter~\ref{ch:deep-bsde} \S\ref{sec:arch-failure-modes}.

\paragraph{(C3) Finite-horizon BSDEJ solver.}  A BSDE-with-jumps solver for the CX model using shared-weight networks and fictitious-play warm-start ([\texttt{solver\_cx\_bsdej\_shared.py}](../../solver_cx_bsdej_shared.py)) achieves spread error $2.6\%$ on the stationary benchmark --- an order of magnitude better than naive per-timestep BSDEJ solvers and essentially the consistency limit for the finite-horizon discretisation. The compensated-martingale Wang et al.\ \cite{wang2023} fix is critical: without compensation the error is $264\%$. Chapter~\ref{ch:deep-bsde} \S\ref{sec:bsdej-solver}.

\paragraph{(C4) Statistical confirmation of tacit collusion.}  Training decentralised MADDPG actor-critic agents on the CX market at $N = 2$ across 20 Hamilton-cluster seeded runs: 15 of 20 (one-sided sign test $p = 0.002$) converge to final spreads strictly above the single-shot Nash value. At $N = 5$, 5 of 5 runs do. The 95\% confidence interval on the mean final spread excludes both Nash and Pareto, placing the outcome in a stable intermediate regime consistent with partial tacit collusion. The mechanism combines pretraining anchor, slow gradient dynamics, value-function generalisation, and decayed exploration. Chapter~\ref{ch:nash} \S\ref{sec:tacit-collusion}.

\section{Relationship to prior work}

The four contributions connect to distinct prior literatures:

\begin{itemize}[leftmargin=2em]
  \item (C1) builds on Cont and Xiong \cite{cont2024} (the exact Algorithm~1 benchmark) and Han--Jentzen--E \cite{han2018} (the deep BSDE template). The neural Bellman variant is simpler than HJE in the stationary regime and delivers better accuracy at lower compute; it is the natural method when time-homogeneity is available.
  \item (C2) builds on Han--Hu--Long \cite{han2022} (MV-FBSDE via fictitious play). The failure-mode analysis has not, to my knowledge, been reported elsewhere --- broadcast-feature-BatchNorm interaction is a known issue in conditional GAN literature \cite{miyato2018}, but the DeepSets symmetric-cancellation in the MV-BSDE context appears to be novel.
  \item (C3) builds on Wang et al.\ \cite{wang2023} (the compensated-martingale fix for jump BSDEs). The shared-weights + warm-start variant is a practical contribution that extends existing BSDEJ solvers to the stationary-CX regime.
  \item (C4) builds on Calvano et al.\ \cite{calvano2020} (tacit collusion in algorithmic pricing via Q-learning) and Cont--Xiong \cite{cont2024} \S6 (the original MADDPG evidence). Our contribution is to reproduce the effect in a continuous-action deep-RL setting with larger samples and a full statistical test.
\end{itemize}

\section{Limitations}
\label{sec:limitations}

Five categories of limitation are worth stating explicitly.

\paragraph{Model stylisation.}  The CX model, while analytically tractable, is a heavy stylisation of real market-making. Real dealer markets involve:
\begin{itemize}
  \item Non-exponential flow-sensitivity (real execution curves are sigmoidal).
  \item Multiple levels of order-book depth (we model only the touch).
  \item Opponent inventory observable to some extent via the order book.
  \item Adverse selection by informed traders (partially addressed in the adverse-selection extension but not the main model).
  \item Fee and rebate structures that can invert the economics of quoting.
\end{itemize}
Our conclusions about Nash, Pareto, and tacit collusion are \emph{qualitative} indicators of market-making behaviour, not specific quantitative predictions about any particular exchange.

\paragraph{Discrete inventory.}  The exact solver works on an integer inventory grid $\{-Q, \ldots, Q\}$, and the neural Bellman solver matches. The diffusion-approximation and continuous-inventory solvers are consistent in the high-arrival-rate limit but systematically differ in the low-rate regime. For very slow markets, the discrete-event dynamics should be solved directly (i.e.\ the jump-BSDE solver).

\paragraph{No common noise.}  The price process in our setup is driftless and independent of inventory in the mean-field limit. Real markets have a common price shock that all agents observe, which introduces a \emph{conditional} mean-field problem: the population law $\mu_t$ depends on the common noise realisation $W^S_t$, and the BSDE becomes conditional on $\mathcal{F}^S_t$. The repo's [\texttt{scripts/common\_noise.py}](../../scripts/common_noise.py) and [\texttt{scripts/conditional\_mv\_common\_noise.py}](../../scripts/conditional_mv_common_noise.py) address this but were not integrated into the main theoretical chain.

\paragraph{MADDPG training instability.}  The 5/20 seeds at $N = 2$ that did \emph{not} finish above Nash suggest substantial training-run-to-run variance. The 95\% CI on the mean spread is consistent with the stable intermediate regime but does not rule out occasional collapses to the Nash equilibrium. Addressing this would require larger sample sizes, variance-reduction techniques (e.g.\ soft target-network updates with smaller $\tau$), or a different actor-critic algorithm less sensitive to exploration schedule.

\paragraph{Proofs of convergence.}  The dissertation is primarily an empirical computational study. Convergence theorems for the neural Bellman solver would require (i) Lipschitz regularity of the Bellman target map (easy), (ii) universal approximation of $V$ by the chosen architecture (standard), (iii) a stochastic-approximation argument for gradient descent to converge to the fixed point (non-trivial). None of these are in the dissertation; see the ``future work'' \S\ref{sec:future-work}.

\section{Future work}
\label{sec:future-work}

Five directions for extension stand out.

\paragraph{Continuous-time BSDEJ with path-sampled jumps.}  The current BSDEJ solvers use Poisson increments \emph{aggregated} over time-discretisation steps $\Delta t$. A more refined scheme would sample actual jump times $\tau_1, \tau_2, \ldots$ explicitly and run the BSDE Euler update piecewise. This reduces time-discretisation bias, especially for large $\Delta t$ or rare jumps, and would tighten the validation error below the current $2.6\%$.

\paragraph{Common noise via deep signatures.}  Hu, Lyons, and collaborators \cite{hu2025} have recently proposed using the \emph{signature} of the common-noise path $(W^S_s)_{s \leq t}$ as a conditioning feature in deep BSDE solvers. This would allow the CX model extension with common noise to be attacked directly. The repo's exploratory scripts in this direction should be extended.

\paragraph{Multi-asset with curse-of-dimensionality regimes.}  The multi-asset extension ($K$ correlated assets, each with its own inventory) has state dimension scaling as $2K$. At $K = 3$ the exact solver requires $10^3$ grid points per dimension and stalls; at $K = 5$ it is infeasible. The deep BSDE methodology scales linearly in $K$ and should be the natural tool; [\texttt{scripts/multiasset\_K5\_CoD.py}](../../scripts/multiasset_K5_CoD.py) explores this regime.

\paragraph{Real LOB data calibration.}  None of this dissertation uses actual market data. A natural follow-up is to calibrate the model parameters $(\lambda^s, \alpha, r, \phi)$ to real LOB snapshots (e.g.\ LOBSTER Nasdaq samples) and compare the predicted equilibrium quotes to observed dealer quotes. This would provide an external validation of the model's descriptive power and a route to empirical tests of the tacit-collusion hypothesis.

\paragraph{Theoretical convergence proofs.}  A rigorous convergence theorem for the neural Bellman solver (or the BSDEJ with shared weights) would be a natural theoretical counterpart to the empirical work here. The argument would likely combine stochastic-approximation theory for the inner gradient descent with Wasserstein contraction for the outer fictitious play.

\section{Closing remarks}

The Cont--Xiong dealer market is a small model: one inventory state, two quotes per side, no information asymmetries in the baseline. Its importance to this dissertation is not that it is realistic but that it is \emph{tractable}: it admits an exact benchmark against which numerical methods can be validated. Once validated, the methods scale to settings where no benchmark exists, and the dissertation's main practical message is that deep BSDE methodology, done carefully, is now good enough to be a reliable component of the applied stochastic-control toolkit.

The less-practical message --- more a meta-lesson --- is in Chapter~\ref{ch:deep-bsde}: architectural decisions that seem innocuous (which normalisation to use, which pooling, which conditioning scheme) can silently destroy the very effect one is trying to study. The failure modes I identified were ones I fell into personally; the preprint and this dissertation are, in large part, the lessons I extracted from debugging my own code. I suspect similar failure modes exist in other deep BSDE applications and that rigorous ablation studies should become standard practice before claiming the absence of an effect.
